% answers_short.tex —— 2015 年·重庆卷·理工农医类 · 简短答案速查
% 由 tools/make_short_answers.py 自 answers.tex 的 \gksolhead/\gkanswer 提取，
% 并与 2015/2015_数学真题_重庆理_答案要点.md 逐题比对后生成，勿手改（改源 MD 后重跑）。

\begin{gknotice}
\textbf{使用说明}：本节只给最终答案便于快速核对，完整推导见解析版。本题库无卷内答案：全部为项目独立求解、经第二人盲解对照（D 级），非官方标准答案。
\end{gknotice}

\gkpart{选择题}
\begin{center}\begin{tabular}{ccccccccccc}
\toprule
题号 & 1 & 2 & 3 & 4 & 5 & 6 & 7 & 8 & 9 & 10 \\
\midrule
答案 & D & B & B & B & A & A & C & C & C & A \\
\bottomrule
\end{tabular}\end{center}

\gkpart{填空题}
\begin{itemize}[leftmargin=2.6em]
  \item[\textbf{11.}] $3$。
  \item[\textbf{12.}] $\dfrac{5}{2}$。
  \item[\textbf{13.}] $\sqrt{6}$。
  \item[\textbf{14.}] $2$。
  \item[\textbf{15.}] $(2,\pi)$。
  \item[\textbf{16.}] $4$ 或 $-6$。
\end{itemize}

\gkpart{解答题}
\begin{itemize}[leftmargin=2.6em]
  \item[\textbf{17.}] (1) $\dfrac{1}{4}$；(2) 分布列为 $P(X=0)=\dfrac{7}{15}$，$P(X=1)=\dfrac{7}{15}$，$P(X=2)=\dfrac{1}{15}$；数学期望 $E(X)=\dfrac{3}{5}$
  \item[\textbf{18.}] (1) 最小正周期 $T=\pi$，最大值 $1-\dfrac{\sqrt{3}}{2}$；(2) $f(x)$ 在 $\left[\dfrac{\pi}{6},\dfrac{5\pi}{12}\right]$ 上单调递增，在 $\left[\dfrac{5\pi}{12},\dfrac{2\pi}{3}\right]$ 上单调递减
  \item[\textbf{19.}] (1) 证明见解析；(2) $\dfrac{\sqrt{3}}{6}$
  \item[\textbf{20.}] (1) $a=0$，切线方程为 $3x-\mathrm{e}y=0$（即 $y=\dfrac{3}{\mathrm{e}}x$）；(2) $a\ge-\dfrac{9}{2}$，即 $a\in\left[-\dfrac{9}{2},+\infty\right)$
  \item[\textbf{21.}] (1) $\dfrac{x^{2}}{4}+y^{2}=1$；(2) $e=\sqrt{6}-\sqrt{3}$
  \item[\textbf{22.}] (1) $a_{n}=3\cdot 2^{n-1}$；(2) 证明见解析
\end{itemize}
